\documentclass[../../main.tex]{subfiles}
\begin{document}
% SOURCE p=257; printed=242; object=ch08-approximate-identity-convergence; number=8.14; region=lower
\begin{theorem}\label{theorem:ch08-approximate-identity-convergence}
设 $\phi\in L^1$ 且 $\int\phi(x)\,\mathrm dx=a$。
\begin{enumerate}
\item\label{item:ch08-approximate-lp}若 $f\in L^p$（$1\leqslant p<\infty$），则当 $t\to0$ 时，$f*\phi_t\to af$ 关于 $L^p$ 范数收敛。
\item\label{item:ch08-approximate-uniform}若 $f$ 有界且一致连续，则当 $t\to0$ 时，$f*\phi_t\to af$ 一致收敛。
\item\label{item:ch08-approximate-local-uniform}若 $f\in L^\infty$，且 $f$ 在开集 $U$ 上连续，则当 $t\to0$ 时，$f*\phi_t\to af$ 在 $U$ 的每个紧子集上一致收敛。
\end{enumerate}
\end{theorem}
% SOURCE p=258; printed=243; object=ch08-approximate-identity-convergence-proof; region=upper-middle
\begin{proof}
令 $y=tz$，有
\begin{align*}
f*\phi_t(x)-af(x)
&=\int[f(x-y)-f(x)]\phi_t(y)\,\mathrm dy\\
&=\int[f(x-tz)-f(x)]\phi(z)\,\mathrm dz
=\int[\tau_{tz}f(x)-f(x)]\phi(z)\,\mathrm dz.
\end{align*}
应用积分形式的 Minkowski 不等式，
\begin{align*}
\|f*\phi_t-af\|_p\leqslant
\int\|\tau_{tz}f-f\|_p|\phi(z)|\,\mathrm dz.
\end{align*}
由\cref{proposition:ch08-lp-translation-continuity}，对每个固定 $z$，$\|\tau_{tz}f-f\|_p\to0$，而它不超过 $2\|f\|_p$，故支配收敛定理给出 (a)。

(b) 的证明相同，只需将 $\|\cdot\|_p$ 换成 $\|\cdot\|_u$。直接取上确界可得
\begin{align*}
\|f*\phi_t-af\|_u\leqslant
\int\|\tau_{tz}f-f\|_u|\phi(z)|\,\mathrm dz.
\end{align*}
一致连续性使被积函数逐点趋零，而 $2\|f\|_u|\phi(z)|$ 是可积支配，故得到一致收敛。

证明 (c)。给定 $\varepsilon>0$，选紧集 $E\subseteq\mathbb R^n$ 使 $\int_{E^c}|\phi|<\varepsilon$。若 $K\subseteq U$ 紧，则 $t$ 足够小时，所有 $x\in K,z\in E$ 均满足 $x-tz\in U$；由紧性及连续性，如\cref{lemma:ch08-compact-uniform-continuity}的论证，
\begin{align*}
\sup_{x\in K,\ z\in E}|f(x-tz)-f(x)|<\varepsilon.
\end{align*}
于是
\begin{align*}
\sup_{x\in K}|f*\phi_t(x)-af(x)|
&\leqslant\sup_{x\in K}
\left[\int_E+\int_{E^c}\right]
|f(x-tz)-f(x)|\,|\phi(z)|\,\mathrm dz\\
&\leqslant\varepsilon\int|\phi|+2\|f\|_\infty\varepsilon.
\end{align*}
由于 $\varepsilon$ 任意，(c) 成立。
\end{proof}
% SUPPLEMENT: id=v14-ch08-approximate-convergence-details; parent=ch08-approximate-identity-convergence; kind=omitted-proof+motivation+hypotheses; origin=independent-derivation
\begin{remark}
\textbf{三种收敛的条件核验（重构新增）。}(a) 中，Young 不等式保证每个 $t>0$ 的卷积作为 $L^p$ 元素存在；差积分的绝对可积性几乎处处也由 Minkowski 定理给出，因为
$\int\|\tau_{tz}f-f\|_p|\phi(z)|\,\mathrm dz\leqslant2\|f\|_p\|\phi\|_1$。平移关于 $L^p$ 范数连续还保证 $z\mapsto\|\tau_{tz}f-f\|_p$ 连续，从而支配收敛的可测性成立。(b) 中有界性保证每点的卷积绝对收敛；$f$ 一致连续使平移差的上确界随平移参数连续，故相应被积函数同样可测。两种情形分别使用可积支配 $2\|f\|_p|\phi|$ 与 $2\|f\|_u|\phi|$，没有要求 $\phi$ 非负。

(c) 中先处理代表问题。记 $M=\|f\|_\infty$。若连续的 $f|_U$ 在某点绝对值大于 $M$，则一个非空开球上的绝对值也大于 $M$，与本质界矛盾。因此 $|f|\leqslant M$ 在 $U$ 上逐点成立。在 $U$ 外的例外零集上将其改成零，既不改变任何卷积值，也保留 $U$ 上原有连续函数，使全局逐点界成立。若 $K=\varnothing$，一致收敛无须证明；否则由紧性及 $U$ 开，可取 $d>0$ 使紧集 $K_d=\{x:\operatorname{dist}(x,K)\leqslant d\}$ 包含于 $U$。$f$ 在 $K_d$ 上一致连续。由于 $E$ 有界，$t\sup_{z\in E}|z|\to0$，故足够小的 $t$ 同时保证 $x-tz\in K_d$ 及所需的统一误差界。$E=\varnothing$ 或只含零点时相应结论直接成立。尾部则由全局界 $2M$ 控制，正是原证最后一项的根据。

动机是把误差写成小平移之差的平均。$\phi\in L^1$ 保证总质量和支配可积；$a$ 可以是任意复数，归一化为一只在希望逼近恒等算子时才需要。有限 $p$ 的限制来自 $L^p$ 平移连续性，不能直接推广到全部 $L^\infty$。一致版本用有界一致连续性替代该条件；局部一致版本用紧子集上的统一连续控制配合全局有界的尾部控制。$U$ 开使紧子集有一个仍在 $U$ 内的厚邻域，而不是只在 $K$ 本身上知道连续性。
\end{remark}
% END SUPPLEMENT: id=v14-ch08-approximate-convergence-details
% SOURCE p=258; printed=243; object=ch08-lebesgue-point-convergence-introduction; region=middle-lower
若对 $\phi$ 施加稍强的条件，还能证明对 $f\in L^p$ 有 $f*\phi_t\to af$ 几乎处处。下面证明把积分按二进区间 $[2^k,2^{k+1}]$ 对应的区域分块并逐块估计；这是 Fourier 分析中的常用方法。
% SOURCE p=258; printed=243; object=ch08-lebesgue-point-convergence; number=8.15; region=lower
\begin{theorem}\label{theorem:ch08-lebesgue-point-convergence}
设对某些 $C,\varepsilon>0$ 有
\begin{align*}
|\phi(x)|\leqslant C(1+|x|)^{-n-\varepsilon},
\end{align*}
并设 $\int\phi(x)\,\mathrm dx=a$（该衰减条件由\cref{corollary:ch02-power-integrability}保证 $\phi\in L^1$）。若 $f\in L^p$、$1\leqslant p\leqslant\infty$，则对 $f$ 的 Lebesgue 集中的每个 $x$，当 $t\to0$ 时均有 $f*\phi_t(x)\to af(x)$。特别地，对几乎所有 $x$，以及 $f$ 连续的每个点，都有此收敛。
\end{theorem}
% SOURCE p=258,259,260; printed=243-245; object=ch08-lebesgue-point-convergence-proof; region=bottom-top
\begin{proof}
若 $x$ 在 $f$ 的 Lebesgue 集中，则对任意 $\delta>0$，存在 $\eta>0$ 使
\begin{align}\label{eq:ch08-lebesgue-point-ball-error}
\int_{|y|<r}|f(x-y)-f(x)|\,\mathrm dy\leqslant\delta r^n,
\qquad r\leqslant\eta.
\end{align}
令
\begin{align*}
I_1&=\int_{|y|<\eta}|f(x-y)-f(x)|\,|\phi_t(y)|\,\mathrm dy,\\
I_2&=\int_{|y|\geqslant\eta}|f(x-y)-f(x)|\,|\phi_t(y)|\,\mathrm dy.
\end{align*}
我们将证明 $I_1\leqslant A\delta$，其中 $A$ 与 $t$ 无关，而 $I_2\to0$。由于 $|f*\phi_t(x)-af(x)|\leqslant I_1+I_2$，这就给出
$\limsup_{t\to0}|f*\phi_t(x)-af(x)|\leqslant A\delta$；再由 $\delta$ 任意即得结论。

估计 $I_1$。当 $\eta/t\geqslant1$ 时取整数 $K$ 使 $2^K\leqslant\eta/t<2^{K+1}$；当 $\eta/t<1$ 时取 $K=0$。将球 $|y|<\eta$ 分为环带
$2^{-k}\eta\leqslant|y|<2^{1-k}\eta$（$1\leqslant k\leqslant K$）及中心球 $|y|<2^{-K}\eta$。在第 $k$ 个环带上，
\begin{align*}
|\phi_t(y)|\leqslant Ct^{-n}\left|\frac yt\right|^{-n-\varepsilon}
\leqslant Ct^{-n}\left(\frac{2^{-k}\eta}t\right)^{-n-\varepsilon},
\end{align*}
而在中心球上 $|\phi_t(y)|\leqslant Ct^{-n}$。因此
\begin{align*}
I_1&\leqslant\sum_{k=1}^KCt^{-n}
\left(\frac{2^{-k}\eta}t\right)^{-n-\varepsilon}
\int_{2^{-k}\eta\leqslant|y|<2^{1-k}\eta}|f(x-y)-f(x)|\,\mathrm dy\\
&\quad+Ct^{-n}\int_{|y|<2^{-K}\eta}|f(x-y)-f(x)|\,\mathrm dy.
\end{align*}
由\cref{eq:ch08-lebesgue-point-ball-error}，在 $\eta/t\geqslant1$ 的情形得到
\begin{align*}
I_1
&\leqslant C\delta\sum_{k=1}^K(2^{1-k}\eta)^nt^{-n}
\left(\frac{2^{-k}\eta}t\right)^{-n-\varepsilon}
+C\delta t^{-n}(2^{-K}\eta)^n\\
&=2^nC\delta\left(\frac\eta t\right)^{-\varepsilon}
\sum_{k=1}^K2^{k\varepsilon}
+C\delta\left(\frac{2^{-K}\eta}t\right)^n\\
&=2^nC\delta\left(\frac\eta t\right)^{-\varepsilon}
\frac{2^{(K+1)\varepsilon}-2^\varepsilon}{2^\varepsilon-1}
+C\delta\left(\frac{2^{-K}\eta}t\right)^n\\
&\leqslant2^nC\left[\frac{2^\varepsilon}{2^\varepsilon-1}+1\right]\delta.
\end{align*}
若 $\eta/t<1$，只有中心球，直接有 $I_1\leqslant C\delta(\eta/t)^n\leqslant C\delta$，也满足最后的统一界。

估计 $I_2$。记 $p'$ 为 $p$ 的共轭指数，$\chi=\chi_{\{|y|\geqslant\eta\}}$。由 Hölder 不等式，
\begin{align*}
I_2&\leqslant\int_{|y|\geqslant\eta}(|f(x-y)|+|f(x)|)|\phi_t(y)|\,\mathrm dy\\
&\leqslant\|f\|_p\|\chi\phi_t\|_{p'}+|f(x)|\|\chi\phi_t\|_1.
\end{align*}
只需证明对每个 $1\leqslant q\leqslant\infty$，$\|\chi\phi_t\|_q\to0$。当 $q=\infty$ 时，
\begin{align*}
\|\chi\phi_t\|_\infty
\leqslant Ct^{-n}(1+\eta/t)^{-n-\varepsilon}
=Ct^\varepsilon(t+\eta)^{-n-\varepsilon}
\leqslant C\eta^{-n-\varepsilon}t^\varepsilon.
\end{align*}
当 $q<\infty$ 时，由换元及\cref{corollary:ch02-radial-integrals}，
\begin{align*}
\|\chi\phi_t\|_q^q
&=\int_{|y|\geqslant\eta}t^{-nq}|\phi(t^{-1}y)|^q\,\mathrm dy
=t^{n(1-q)}\int_{|z|\geqslant\eta/t}|\phi(z)|^q\,\mathrm dz\\
&\leqslant C_1t^{n(1-q)}\int_{\eta/t}^\infty r^{n-1-(n+\varepsilon)q}\,\mathrm dr\\
&=C_2t^{n(1-q)}(\eta/t)^{n-(n+\varepsilon)q}
=C_3t^{\varepsilon q}.
\end{align*}
这里 $C_1,C_2,C_3$ 可以依赖固定的 $q,\eta,n,C,\varepsilon$，但与 $t$ 无关。因此每个 $q$ 的范数均为 $O(t^\varepsilon)$，故 $I_2\to0$，证明完成。
\end{proof}
% SUPPLEMENT: id=v14-ch08-lebesgue-point-details; parent=ch08-lebesgue-point-convergence; kind=omitted-proof+motivation+hypotheses; origin=independent-derivation
\begin{remark}
\textbf{端点、常数及逐点结论（重构新增）。}衰减假设首先使 $\phi$ 有界可积，继而 $\phi\in L^q$ 对所有 $1\leqslant q\leqslant\infty$ 成立：有限 $q$ 时 $\int|\phi|^q\leqslant\|\phi\|_\infty^{q-1}\|\phi\|_1$。伸缩后仍如此。因此对任意 $f\in L^p$，用共轭指数的卷积命题可知 $f*\phi_t(x)$ 在\emph{每个} $x$ 存在，不会在一个随 $t$ 变化的例外集上失去定义。

$f\in L^p$ 蕴含局部可积：有限测度紧集上用 Hölder（$p=1$ 直接，$p=\infty$ 用有界性）。Lebesgue 点定义中的球体积为 $m(B_1)r^n$，把这个固定常数吸收到任意小参数中，即得\cref{eq:ch08-lebesgue-point-ball-error}。环带的外半径至多 $\eta$，所以每次使用该估计均在合法范围。几何级数的分母 $2^\varepsilon-1$ 为正，正是 $\varepsilon>0$ 的用途。$\eta/t\geqslant1$ 时有 $2^{-K}\eta/t<2$ 及 $(\eta/t)^{-\varepsilon}2^{(K+1)\varepsilon}\leqslant2^\varepsilon$，给出与 $\eta,t,\delta$ 均无关的 $A$。原文统一使用 $2^K\leqslant\eta/t$ 的一句在 $\eta/t<1$ 时不适用；这里已把只有中心球的情形直接算完。

尾部积分的幂满足 $(n+\varepsilon)q>n$，因此径向积分收敛；其尺度指数为
$n(1-q)-n+(n+\varepsilon)q=\varepsilon q$。$p=1$ 用 $p'=\infty$ 的上确界估计，$p=\infty$ 用 $p'=1$ 的尾积分估计，两端均已覆盖。取极限时先固定 $\delta$ 及相应 $\eta$，令 $t\downarrow0$ 消去尾项，再令 $\delta\downarrow0$；不能反过来让依赖 $\eta$ 的尾常数被当成统一常数。

由 Lebesgue 微分定理，几乎每点为 Lebesgue 点；若 $f$ 在一点连续，则该点附近的平均振幅不超过相应邻域上的最大振幅，并趋零，所以该点也是 Lebesgue 点。动机是以局部平均的小误差控制中心部分，以核的严格超维数衰减控制远处部分。结论中的“每个 Lebesgue 点”比只说几乎处处更精确；不要求 $\phi$ 非负或积分为一，极限自然是 $af(x)$。
\end{remark}
% END SUPPLEMENT: id=v14-ch08-lebesgue-point-details
\subfile{section-02-cutoffs.tex}
\end{document}
